\section{Discussion}
\td{check revised version below.}
\subsection{Principal findings}
Within the same discharge summaries, stigmatizing language was about eight times less frequent in the instructions written for patients than in the narrative written for clinicians. Clinicians therefore already moderate their language when they expect the patient to read it. This moderation was not uniform. For patients with substance use disorder, serious mental illness, or sickle cell disease, the groups for whom stigmatizing documentation has been most consistently reported,\cite{himmelstein2022,goddu2018} discharge instructions contained three to four times the average rate of stigmatizing language, and the gap between clinician-facing and patient-facing text was narrower. The patients most exposed to stigma in clinical documentation were thus also the least protected from it in the text written for them.

\subsection{Audience shapes documentation}
Stigmatizing language in discharge summaries is concentrated in the sections clinicians write for each other. Previous studies analyzed the clinical note as a single unit and therefore could not make this distinction.\cite{himmelstein2022,harrigian2023,xue2026,walker2025,scroggins2025} Our within-note design compared sections written by the same clinician about the same patient and admission, so the difference between sections cannot be explained by who the patient is or who wrote the note. What differs is the intended reader.

The keyword analysis shows that this difference has two sources. The first is content: clinicians used stigma-related keywords 6.55 times as often per word in BHC as in DI. This is expected given the purpose of each section. BHC and HPI recount what happened during the admission, including adherence, behavior, and interactions with staff, whereas DI tells the patient what to do next. The second is framing: when a keyword did appear, it was 1.54 times as likely to be used in a stigmatizing way in BHC as in DI. Writing for patients, clinicians thus both raise these topics less often and phrase them more neutrally when they do.

This finding has a practical consequence. Clinicians evidently write more carefully when they expect the patient to read the text. Under the 21st Century Cures Act, however, patients can now access most clinical notes, including the clinician-facing sections.\cite{onc2020cures} The care taken in DI does not yet extend to the sections that contain most of the stigmatizing language, even though patients can now read them as well.

\subsection{Where the gap narrows}
The cohorts with the highest DI rates were those with the narrowest gaps, and their BHC rates were also above average. Stigmatizing language in these patients' instructions therefore does not reflect cleaner clinician-facing text but language that was carried into the DI. One explanation is that discharge instructions for these conditions must address the very topics the lexicon captures, such as substance use, pain medication, and adherence, so stigma-related terms are harder to avoid. Some of this language may be clinically necessary. The concern is how it is framed. Obesity showed the opposite pattern: its DI rate was the lowest of all cohorts while its BHC rate was average, consistent with instructions focused on routine management. The gap also narrowed with age. Whether this reflects different instruction content for older patients or instructions written for caregivers cannot be determined from these data.

These findings also show why relative measures should be reported alongside absolute rates. A narrower gap can result from less stigmatizing language in clinician-facing text or from more in patient-facing text. Here it reflected the latter, which the interaction RRs alone would not have revealed.

\subsection{Implications}
First, efforts to reduce stigmatizing language in patient-facing text can be targeted. Review or writing support at the point of discharge could focus on patients with substance use disorder, serious mental illness, and sickle cell disease, for whom the burden is concentrated, rather than on all notes. Second, the classifier provides a scalable way to audit documentation, for example to monitor language by section, service, or patient group over time. Third, as large language models are increasingly used to draft discharge instructions and patient-friendly summaries,\cite{zaretsky2024generative} they risk reproducing the language of the clinician-facing narrative from which they are generated. Section-aware evaluation of model outputs, using the approach described here, can determine whether such tools reduce or propagate stigmatizing language.
